% Part: first-order-logic
% Chapter: beyond
% Section: many-sorted-logic

\documentclass[../../../include/open-logic-section]{subfiles}

\begin{document}

\olfileid{fol}{byd}{msl}

\olsection{د ډېرو ډولونو منطق}

د لومړۍ درجې په منطق کښې متغيرونه او کميت ټاکونکي پر يوه
!!{domain} ګرځي۔ خو ډېر ځله د څو (يو له بله بېلو) !!{domain}s درلودل
ګټور وي: د بېلګې په توګه، ښايي د عددونو !!a{domain}، د هندسي څيزونو
!!a{domain}، له عددونو څخه عددونو ته د تابعو !!a{domain}، د ابېلي
ګروپونو !!a{domain} او داسې نور وغواړئ۔

د ډېرو ډولونو منطق دا ډول چوکاټ برابروي۔ پيل د ``ډولونو'' له يوه لړ
څخه کېږي---د يوه څيز ``ډول'' هغه ``!!{domain}'' ښيي چې څيز بايد پکې
اوسېږي۔ بيا د هر ډول دپاره !!{variable}s او کميت ټاکونکي، او (په عام
ډول) د هر ډول دپاره يو !!{identity} لرو۔ تابعې او اړيکې هم د هغو
څيزونو د ډولونو له مخې ``ټايپ'' کېږي چې د ارګومېنټونو په توګه ئې
اخيستلے شي۔ له دې پرته د لومړۍ درجې د منطق عادي قاعدې ساتل کېږي، او
د کميت ټاکونکو د قاعدو بڼې د هر ډول دپاره تکرارېږي۔

د بېلګې په توګه، د نړيوالو اړيکو د مطالعې دپاره ښايي داسې ژبه وټاکو
چې دوه ډوله څيزونه، فرانسوي اتباع او جرمن اتباع، ولري۔ ښايي د لومړي
ډول د څيزونو دپاره يوه يوځاييزه اړيکه، ``شراب څښي''، د دويم ډول د
څيزونو دپاره بله يوځاييزه اړيکه، ``ورسټ خوري''، او يوه دوه ځاييزه
اړيکه، ``څو ملتيزه واده شوې جوړه جوړوي''، ولرو۔ وروستۍ اړيکه دوه
ارګومېنټونه اخلي: لومړے د لومړي ډول او دويم د دويم ډول وي۔ که $a$،
$b$، $c$ متغيرونه پر فرانسوي اتباعو او $x$، $y$، $z$ متغيرونه پر
جرمن اتباعو وګرځوو، نو
\[
\lforall[a][\lforall[x][(\Atom{\Obj{MarriedTo}}{a,x} \lif
(\Atom{\Obj{DrinksWine}}{a} \lor \lnot \Atom{\Obj{EatsWurst}}{x})]]
\]
دا وايي چې که کوم فرانسوے له يوه جرمن سره وادۀ شوے وي، نو يا هغه
فرانسوے شراب څښي يا هغه جرمن ورسټ نۀ خوري۔
\textbf{د ژباړې د نسخې سمون (OLFOL-030):} سرچينې د دويم عمومي کميت
ټاکونکي د متغير پرانيستونکې نښه غورځولې وه؛ دلته هغه نښه بېرته
ايښوول شوې ده۔

د ډېرو ډولونو منطق په طبيعي ډول د لومړۍ درجې په منطق کښې ځايېدلے شي:
د ډېرو ډولونو د !!{domain}s ټول څيزونه په يوه د لومړۍ درجې !!{domain}
کښې يوځاے کوو، د ډولونو د پېژندلو دپاره يوځاييز !!{predicate}s کاروو
او کميت ټاکونکي نسبي کوو۔ د بېلګې په توګه، د پاسنۍ بېلګې اړونده د
لومړۍ درجې ژبه به د نورو بيان شويو اړيکو تر څنګ ``$\Obj{German}$''
او ``$\Obj{French}$'' يوځاييز !!{predicate}s ولري، او د ډول شرطونه
به ترې لرې وي۔ يو ډول‌لرونکے کميت ټاکونکے $\lforall[x][!A]$، چې $x$
د جرمن ډول !!a{variable} وي، داسې ژباړل کېږي:
\[
\lforall[x][(\Atom{\Obj{German}}{x} \lif !A)].
\]
بايد داسې بديهي اصول هم ورزيات کړو چې ډولونه بېل تضمين کړي---لکه
$\lforall[x][\lnot (\Atom{\Obj{German}}{x} \land
  \Atom{\Obj{French}}{x})]$---او هم داسې بديهي اصول چې تضمين کړي
``شراب څښي'' يوازې د $\Atom{\Obj{French}}{x}$ پريديکات پوره کوونکو
څيزونو دپاره صدق کوي، او داسې نور۔ له دې تړونونو او بديهي اصولو سره
دا ښودل ګران نۀ دي چې د ډېرو ډولونو !!{sentence}s د لومړۍ درجې
!!{sentence}s ته، او د ډېرو ډولونو !!{derivation}s د لومړۍ درجې
!!{derivation}s ته ژباړل کېږي۔ همدارنګه، د ډېرو ډولونو !!{structure}s
د لومړۍ درجې اړوندو !!{structure}s ته او برعکس ``ژباړل'' کېږي؛ نو د
ډېرو ډولونو د منطق دپاره هم د بشپړتيا قضيه لرو۔

\end{document}
